% 04_results_A.tex — Results A: the neuroscience battery, quantified (P2-E8-5).
% Maier voice per ../STYLE.md. Every number traces to a committed §R record
% (tools/xai_study/phaseA_kording/out/*.json) via the E6-1 leaderboard and Fig. 3.
% Honesty contract (plan.md): A is the low-faithfulness calibration baseline, not the
% headline; we report what the methods recover against the KNOWN mechanism, and we do
% NOT claim any of them recovered meaning or algorithm.

\section{The neuroscience battery, quantified}\label{sec:results-a}

We begin where the field's doubt began. Jonas and Kording asked whether a
neuroscientist, armed with the standard analysis toolkit, could understand a
microprocessor, and answered that the toolkit yields rich, suggestive structure
that nonetheless misses the mechanism~\cite{jonas2017could}. Their study could
diagnose the failure but not grade it: on the 6502 they too lacked a number for
how wrong each method was, because they had no causal ground truth to score
against. We supply that number. Running the same battery on the bit-exact,
fully known VCS---where the true causal role of every register and RAM cell is
recoverable by exact intervention (\S\ref{sec:methods})---turns Kording's
qualitative lesson into a measured one. The result is the calibration baseline
for the rest of the paper: classical analyses find structure and score
\emph{low} in faithfulness against the known register-transfer account, while
the oracle-as-method positive control sits at $1.0$ by construction
(Fig.~\ref{fig:phaseA}).

The headline is a spread, not a single value. Across the eight analyses A1--A8,
faithfulness against the oracle ranges from $0.116$ to $1.0$, averaging
$0.49 \pm 0.24$ over the phase~\cite{leaderboardP2}. Read against the perfect
intervention oracle that anchors the same axis, the typical classical method is
closer to chance than to the mechanism. The structure these methods report is
real---the data-flow graphs, the tuning, the correlations all exist in the
state tensor---yet structure is not the same as cause, and the gap between the
two is what the score measures. We walk the battery once, naming each method,
what it returns, and how far that lands from the truth the substrate already
knows.

\paragraph{Connectomics and lesions: the wiring is largely recoverable.}
A1 reconstructs the dependency graph over state variables---a connectome of
read/write edges---and scores it against the \emph{true} data-flow graph by
graph-matched precision and recall. It reaches a faithfulness of $0.41$
(mean data-flow $F_1$ vs.\ the oracle)~\cite{leaderboardP2}, the structure
half-right and half-spurious. A2, the exhaustive single-unit lesion, fares
best of the classical set: clamping each RAM cell in turn and ranking it by
the behavioral damage it does yields an importance map whose Spearman
correlation with the unit's true causal role is $0.99$ (mean $0.987$ over the
six core games)~\cite{leaderboardP2}. Lesioning, in short, is causal, and it
shows: A2 is the one classical method whose number sits near the ceiling.

\paragraph{Yet a near-perfect rank order hides an interaction blind-spot.}
The $\rho = 0.99$ is seductive, and here the honest reading matters. A
single-unit lesion changes one cell at a time, so it is blind to causes that
act only in concert. On Space Invaders, the per-unit ranking misses $25\%$ of
the interactions the oracle records, and a single super-additive pair produces
a behavioral swing of up to $1044$ screen-pixels of change beyond the sum of
its parts~\cite{leaderboardP2}. The map of \emph{which units matter} is right;
the model of \emph{how they combine} is absent. Hence even the strongest
classical method recovers a ranking, not the computation---an early instance of
the theme the paper makes precise later: a faithful summary of effects is not
the algorithm that produces them.

\paragraph{Tuning curves: structure that points the wrong way.}
A3 fits per-unit tuning to a game variable, the workhorse of systems
neuroscience. On the VCS it is actively misleading. We measure a
\emph{spurious-tuning rate}---the fraction of strongly-tuned units whose tuning
does not match their true causal role---and on Pong it is $1.0$: every unit
that looks tuned to a game concept is, by the oracle, not causally that
variable at all~\cite{leaderboardP2}. The method's faithfulness collapses to
$0.12$, the lowest in the battery. Tuning is correlational; a cell can track a
variable it does not control, and on a machine where many cells co-vary with
the beam clock, apparent tuning is cheap. This is the tuning-curve trap stated
in numbers, and it foreshadows the present-versus-used distinction we return to
for linear probing.

\paragraph{Correlation structure and field potentials: real, and largely epiphenomenal.}
A4 reads the pairwise and global correlation structure of the state---the
spike-word analysis---and reproduces the familiar weak-pairwise/strong-global
signature, but against the true coupling its faithfulness is only $0.29$, with
sufficiency near zero ($S = 0.018$)~\cite{leaderboardP2}: the correlations are
genuine yet predict almost nothing about held-out interventions. A5 pools
regional activity into power spectra, the local-field-potential analogue, and
scores $0.40$; much of the variance it captures is the known frame and scanline
clocks~\cite{leaderboardP2}---signal that is real, periodic, and causally
epiphenomenal. Both methods see the machine breathing without seeing it
compute.

\paragraph{Granger causality: temporal precedence is not causation.}
A6 infers directed CPU/TIA/RIOT edges by Granger causality and scores its
edges against the true data-flow. On Pong its faithfulness is $0.0$: the
false-edge and missed-edge rates are both $1.0$~\cite{leaderboardP2}. The
reason is structural and instructive. The VCS is a clock-locked machine in
which nearly everything is statistically predictable from everything else one
cycle earlier, so a method that equates predictive precedence with causation
infers spurious bidirectional edges everywhere. Granger asks the wrong
question of a synchronous substrate, and the exact answer the substrate can
give---intervene and re-run---exposes it completely. \emph{Note that} this is
not a tuning failure we could repair with a longer lag or a stricter
threshold; precedence and causation simply come apart on this machine, and only
the intervention oracle can tell them apart.

\paragraph{Dimensionality reduction: NMF edges out PCA, both mediocre.}
A7 decomposes the state tensor into latent components and matches them to the
known signals (clock, read/write, vsync). Non-negative matrix factorization
recovers a matched-component fraction of $0.60$ (best $0.80$ on a single game),
slightly ahead of PCA at $0.60$ on average, the non-negativity of NMF a better
fit to a substrate of non-negative activations~\cite{leaderboardP2}. The
ordering is the expected one, and the magnitude is the point: neither method
clears two-thirds of the known components, so the latent space is interpretable
in part and opaque in the rest. A7's headline faithfulness of $0.60$ makes it
the second-strongest classical analysis, behind only the causal lesion.

\paragraph{Whole-state recording: a perfect record is not an explanation.}
A8 simply records the full RAM and register state over time---the descriptive
baseline, the complete observation no real experiment ever has. It is faithful
and sufficient by definition ($F = 1.0$, $S = 1.0$), and yet its minimality
score is $0.07$~\cite{leaderboardP2}: a recording of everything is, at the
algorithmic level, an explanation of nothing. A8 is the limiting case that
makes the triad concrete. Faithfulness and sufficiency without minimality is a
log file, not an account of the computation---the first hint that the three
axes can pull apart, which the semantic-gap result later turns into the paper's
centre of gravity.

\paragraph{The quantified Kording.} Taken together, the battery confirms the
expected, important result and puts a number on it. Classical neuroscience
methods, run on a system whose answer is known exactly, find rich structure and
score low in faithfulness---$0.49$ on average, with only the causal lesion (A2,
$\rho = 0.99$) and dimensionality reduction (A7, $0.60$) rising clearly above
chance, and Granger causality (A6) sitting at zero. The advance over Jonas and
Kording is threefold: the subject is the \emph{architectural} register-transfer
level rather than the transistor; the score is a \emph{quantitative} $F$ against
an exact oracle rather than a qualitative verdict; and the contrast is sharp,
because the causal operators we turn on the same machine in
Sections~\ref{sec:results-b}--\ref{sec:results-c} reach the ceiling that the
classical battery cannot. \emph{The gap is therefore the method, not the
machine.} We mean this strictly as a calibration claim: A grades the
faithfulness floor, and even its best members recover a ranking, a graph, or a
record---never the algorithm that the program runs. That residual is the
problem the rest of the paper measures.

\begin{figure}[t]
  \centering
  \includegraphics[width=\textwidth]{fig3_phaseA_battery.pdf}
  \caption{\textbf{The Kording neuroscience battery, scored against the exact
  oracle.} Each analysis A1--A8 is run on the bit-exact, fully known VCS and
  graded for faithfulness against the \S\ref{sec:methods} intervention oracle;
  bars are coloured by method tradition, and the oracle-as-method positive
  control is drawn at faithfulness $=1.0$ (dashed reference). The classical
  battery averages $0.49$, with only the causal single-unit lesion (A2) and
  dimensionality reduction (A7) clearly above chance. \emph{Note that} the four
  callouts isolate the mechanism behind the low scores: A6 Granger causality
  reaches $F=0$ on Pong (temporal precedence read as causation on a clock-locked
  machine), A3 tuning shows a spurious-tuning rate of $1.0$ on Pong
  (strongly-tuned units whose tuning is not their true causal role), A2 attains
  a single-unit importance correlation of $\rho=0.99$ yet misses $25\%$ of
  interactions with a super-additive pair swinging up to $1044$ pixels, and A7
  NMF edges out PCA on matched components while both stay below two-thirds. The
  figure is the calibration baseline that the causal contrast of
  Fig.~\ref{fig:faithfulness} sharpens. Every plotted value reads from a
  committed per-game record; no number is re-run or estimated.}
  \label{fig:phaseA}
\end{figure}

% BIBTEX-NEEDED: full entries for keys cited above that are NOT already in
% references.bib. (jonas2017could and sundararajan2017axiomatic already exist.)
% Please merge the following into references.bib at SM Review.
%
% @misc{leaderboardP2,
%   author       = {{Paper~2 authors}},
%   title        = {Cross-tradition faithfulness--plausibility leaderboard for the
%                   differentiable {Atari} {VCS} interpretability audit (Phases A--C)},
%   year         = {2026},
%   howpublished = {Committed analysis record,
%                   \texttt{tools/xai\_study/compare/out/leaderboard.json};
%                   aggregates the per-game §R records under
%                   \texttt{tools/xai\_study/phaseA\_kording/out/}},
%   note         = {Internal data artifact; 257 per-game records aggregated into 31
%                   method rows including the oracle positive control. Generated by
%                   \texttt{compare/leaderboard.py}. Accessed 2026-06-23.}
% }
